% Part: history
% Chapter: set-theory
% Section: cantor-plane

\documentclass[../../../include/open-logic-section]{subfiles}

\begin{document}

\olfileid{his}{set}{cantorplane}
\olsection{রেখা ও তল নিয়ে কান্টর}

শ্র্যোডার--বার্নস্টাইন উপপাদ্যের প্রমাণকে ঘিরে কিছু ঘটনা
\olref[his][set][pathology]{sec}-এ আলোচিত ইতিহাসের সঙ্গে
জড়িত। মনে করো, ১৮৭৭ সালে কান্টর প্রমাণ করেছিলেন যে
একটি বর্গক্ষেত্রে যত বিন্দু আছে, তার একটি বাহুতেও ঠিক তত
বিন্দু আছে। এখানে তাঁর প্রথম প্রমাণচেষ্টাটি দেখব।

$\unitline$ বলতে একক রেখাখণ্ড, অর্থাৎ $[0,1]$-এর
বিন্দুগুলির সেট বোঝাই। আর $\unitsquare$ বলতে একক
বর্গক্ষেত্র, অর্থাৎ $\unitline \times \unitline$-এর
বিন্দুগুলির সেট বোঝাই। এই ভাষায় কান্টর প্রমাণ করেছিলেন,
$\cardeq{\unitline}{\unitsquare}$। ডেডেকিন্ডকে লেখা তাঁর
একটি চিঠিতে মূলত নিচের যুক্তিটি ছিল।

\begin{thm}\ollabel{thm:cantorplane}
$\cardeq{\unitline}{\unitsquare}$
\end{thm}

\begin{proof}[প্রমাণ: প্রথম অংশ]
যেকোনো $a, b \in \unitline$ স্থির করো। তাদের দ্বিমিক
বিস্তারের অঙ্কগুলি যথাক্রমে $a_1$, $a_2$, \dots,
এবং $b_1$, $b_2$, \dots।
% BN-SRC-489: specify a canonical binary representation, including the endpoint 1;
% the frozen witness 0.1010... is f(1,0) under this valid convention.
এখানে একের চেয়ে ছোট সংখ্যার জন্য শেষ পর্যন্ত শুধু
এক চলতে থাকে না, এমন বিস্তারটি বেছে নিই; তাই দুটি
বিস্তার থাকলে শূন্যে শেষ হওয়া রূপটি নিই। একের জন্য
নিই সব অঙ্ক একযুক্ত বিস্তার। তাহলে:
\begin{align*}
a &= 0.a_1a_2a_3a_4\dots\\
b &= 0.b_1b_2b_3b_4\dots
\intertext{এবার $f \colon \unitsquare \to \unitline$ অপেক্ষকটি নিই, যেখানে}
f(a, b) & = 0.a_1b_1a_2b_2a_3b_3a_4b_4\dots
\end{align*}
এটি একটি !!a{injection}। প্রকৃতপক্ষে, এভাবে পাওয়া
অঙ্কের ধারা শেষ পর্যন্ত শুধু এক হতে পারে কেবল দুই
ইনপুটই এক হলে; তখন ফলও এক। অন্য সব ক্ষেত্রে,
ফলের আরেকটি দ্বিমিক বিস্তার থাকলেও তা নির্বাচিত
ইনপুটের অঙ্কগুলি আন্তঃবুনে পাওয়া যায় না। অতএব
$f(a, b) = f(c,d)$ হলে সব
$n \in \Nat$-এর জন্য $a_n = c_n$ এবং $b_n = d_n$;
সুতরাং $a = c$ ও $b = d$।
\end{proof}

কিন্তু ডেডেকিন্ড কান্টরকে দেখিয়েছিলেন, এতটুকুতে
মূল প্রশ্নের উত্তর মেলে না। উদাহরণ হিসেবে
% BN-SRC-489: valid non-surjectivity witness for the canonical convention.
$0.00\dot{1}\dot{0} = 0.0010101010\ldots$
বিবেচনা করো। এটি $f(a,b)$ হলে অঙ্কগুলি আলাদা
করে পড়ে পেতে হয়:
\begin{align*}
a&= 0.011111\ldots = \frac{1}{2}\\
b&= 0
\end{align*}
কিন্তু নির্ধারিত নিয়মে অর্ধেকের দ্বিমিক বিস্তার
$0.1000\ldots$; $0.011111\ldots$ রূপটি আমরা
ইনপুটের জন্য নিইনি। আর বিবেচিত ফলটির নিজের
বিকল্প দ্বিমিক বিস্তার নেই। তাই এই বিন্দুটি
$f$-এর কোনো ইনপুট থেকে পাওয়া যায় না।

% BN-SRC-490: the frozen source calls these recurring *decimal* expansions,
% although the construction and dual-representation issue are binary.
সারকথা, কিছু সংখ্যার দ্বিমিক বিস্তারের দুটি রূপ
থাকার কারণে কান্টরের $f$ অপেক্ষকটি
!!a{injection}, কিন্তু !!a{surjection} নয়।
সুতরাং তিনি এতক্ষণে শুধু
$\cardle{\unitsquare}{\unitline}$ দেখিয়েছেন,
$\cardeq{\unitsquare}{\unitline}$ দেখাননি।

কান্টর প্রায় সঙ্গে সঙ্গেই ডেডেকিন্ডকে উত্তর দেন।
মূলত তিনি নিচের মতো করে প্রমাণটি সম্পূর্ণ
করার ইঙ্গিত দেন:

\begin{proof}[প্রমাণ: সমাপ্ত অংশ]
আমরা দেখিয়েছি,
$\cardle{\unitsquare}{\unitline}$।
অন্যদিকে $\unitline$ থেকে $\unitsquare$-এ
একটি !!a{injection} স্পষ্টই আছে:
রেখাখণ্ডটিকে বর্গক্ষেত্রের একটি বাহু বরাবর
বসালেই হয়। তাই
$\cardle{\unitline}{\unitsquare}$ এবং
$\cardle{\unitsquare}{\unitline}$।
শ্র্যোডার--বার্নস্টাইন উপপাদ্য
(\olref[sfr][siz][sb]{thm:schroder-bernstein})
থেকে $\cardeq{\unitline}{\unitsquare}$।
\end{proof}

অবশ্য কান্টর নিজে শেষ পঙ্‌ক্তিটি এই ভাষায়
লিখে প্রমাণ শেষ করতে পারতেন না: তখনও
শ্র্যোডার--বার্নস্টাইন উপপাদ্য প্রমাণিত হয়নি।
পরে কান্টর এই বক্তব্যকে সাধারণ অনুমান
হিসেবে প্রকাশ করলেও ১৮৯৭ সালের আগে
তার সন্তোষজনক প্রমাণ পাওয়া যায়নি।
(তাই ১৮৭৭ সালের পরের দিকে কান্টর
\olref{thm:cantorplane}-এর অন্য একটি
প্রমাণ দিয়েছিলেন, যাতে শ্র্যোডার--বার্নস্টাইন
উপপাদ্য ব্যবহার করা হয়নি।)

\end{document}
